\chapter{Fehlerrechnung}

\setcounter{equation}{0}

Für die Begründung der folgenden Formeln siehe das Anleitungsdokument zum physikalischen Anfängerpraktikum an der Universität Heidelberg, im Teil zur Fehlerrechnung.

\section{Fehlerfortpflanzung}

Im Folgenden werden Funktionen definiert. Definitions- und Wertebereiche setzen sich stets aus den reellen Zahlen $\R$ zusammen.

\begin{enumerate}
    \item   $f(x_1, \dots, x_n)$ Eine beliebige Funktion in n Variablen.
    \item   $f_1(x_1, \dots, x_n) := x_1 + \dots + x_n$
    \item   $f_2(x_1, \dots, x_n) := x_1^{a_1} \cdot \dots \cdot x_n^{a_n}$ \hfill $a_1, \dots, a_n \in \R$
    \item   $f_3(x,y) := x + y$
    \item   $f_4(x) := a x$  \hfill $a \in \R$
    \item   $f_5(x_1, \dots, x_n) := a x_1 \cdot \dots \cdot x_n$  \hfill $a \in \R$
    \item   $f_6(x_1, \dots, x_n, y_1, \dots, y_m) = \displaystyle\frac{x_1^{a_1} \cdot \dots \cdot x_n^{a_n}}{y_1^{b_1} \cdot \dots \cdot y_m^{b_m}}$ \hfill $a_1, \dots, a_n, b_1, \dots, b_m \in \R$
\end{enumerate}

Das allgemeine Fehlerfortpflanzungsgesetz für f lautet:
\begin{equation}
    \Delta f(x_1, \dots, x_n) = \sqrt{\bigg(\frac{\partial f}{\partial x_1} \cdot \Delta x_1\bigg)^2 + \dots + \bigg(\frac{\partial f}{\partial x_n} \cdot \Delta x_n\bigg)^2}    \label{eq:fehler_allg}
\end{equation}

Für $f_1$ bleibt dies unverändert.

Aus Gleichung \eqref{eq:fehler_allg} folgt für $f_2$ im relativen Fehler:
\begin{equation}
    \frac{\Delta f_2(x_1, \dots, x_n)}{f_2(x_1, \dots, x_n)} = \sqrt{\bigg(a_1 \frac{\Delta x_1}{x_1}\bigg)^2 + \dots + \bigg(a_n\frac{\Delta x_n}{x_n}\bigg)^2} \label{eq:fehler_mult_allg}
\end{equation}

Aus Gleichung \eqref{eq:fehler_allg} folgt für $f_3$:
\begin{equation}
    \Delta f_3(x,y) = \sqrt{(\Delta x)^2 + (\Delta y)^2} \label{eq:fehler_add_einfach}
\end{equation}

Aus Gleichung \eqref{eq:fehler_allg} folgt für $f_4$:
\begin{equation}
    \Delta f_4(x) = a \cdot \Delta x \label{eq:fehler_const_einfach}
\end{equation}

Aus den Gleichungen \eqref{eq:fehler_allg} und \eqref{eq:fehler_mult_allg} folgt für $f_5$ im relativen Fehler:
\begin{equation}
    \frac{\Delta f_5(x_1, \dots, x_n)}{f_5(x_1, \dots, x_n)} = \sqrt{\bigg(\frac{\Delta x_1}{x_1}\bigg)^2 + \dots + \bigg(\frac{\Delta x_n}{x_n}\bigg)^2} \label{eq:fehler_const_mult}
\end{equation}

Aus den Gleichung \eqref{eq:fehler_allg} und \eqref{eq:fehler_mult_allg} folgt für $f_6$ im relativen Fehler:
\begin{equation}
    \frac{\Delta f_6}{f_6}  = \sqrt{\bigg(a_1\frac{\Delta x_1}{x_1}\bigg)^2 + \dots + \bigg(a_n\frac{\Delta x_n}{x_n}\bigg)^2 + \bigg(b_1\frac{\Delta y_1}{y_1}\bigg)^2 + \dots + \bigg(b_m\frac{\Delta y_m}{y_m}\bigg)^2} \label{eq:fehler_div}
\end{equation}

\section{Konfidenzintervalle}

In diesem Praktikum wird die folgende Gleichung verwendet, um zwei (eventuell fehlerbehaftete) Werte miteinander zu vergleichen. Die Zahl n gibt die Anzahl an Konfidenzintervallen ($\sigma$-Umgebungen) um einen Literaturwert an, innerhalb derer sich ein Messwert befindet.
\begin{equation}
    n = \frac{|x_{mess} - x_{lit}|}{\sqrt{(\Delta x_{mess})^2 + (\Delta x_{lit})^2}} \label{eq:fehler_konf}
\end{equation}

\section{Mittelwert und Standardabweichung}

Sei X eine zufallsverteilte Variable mit einem Wertebereich der Mächtigkeit N. Sei $\{x_1, \dots, x_n\}$ dieser Wertebereich.

Die Berechnung des Mittelwertes von X erfolgt dann nach der Gleichung:
\begin{equation}
    \overline{x} = \frac{1}{N} \sum_{i=1}^N x_i \label{eq:mittelwert}
\end{equation}

Die Berechnung der Standardabweichung von X erfolgt nach der Gleichung:
\begin{equation}
    S_E' = \sqrt{\frac{1}{N} \sum_{i=1}^N (\overline{x} - x_i)^2} \label{eq:standardabweichung}
\end{equation}

Die Berechnung des Fehlers einer Einzelmessung von X erfolgt nach der Gleichung:
\begin{equation}
    S_E = \sqrt{\frac{1}{N - 1} \sum_{i=1}^N (\overline{x} - x_i)^2} \label{eq:fehler_einzelmessung}
\end{equation}

Die Berechnung des Fehlers des Mittelwertes von X erfolgt nach der Gleichung:
\begin{equation}
    S_M = \sqrt{\frac{1}{N(N - 1)} \sum_{i=1}^N (\overline{x} - x_i)^2}~\label{eq:fehler_mittelwert}
\end{equation} 

Nun können die einzelnen $x_i$ auch fehlerbehaftet sein. Seien $\Delta x_i$ die Fehler der einzelnen Messwerte. Dann gilt:
\begin{equation}
    \Delta \overline{x} = \sqrt{\frac 1 N \cdot S_E^2 + \frac{1}{N^2} \cdot \sum_{i=1}^N (\Delta x_i)^2} \label{eq:fehler_messreihe_allg}
\end{equation}

Sind die Fehler aller $x_i$ gleich, $\Delta x_i =: \Delta x$, so vereinfacht sich dies zu:
\begin{equation}
    \Delta \overline{x} = \sqrt{\frac 1 N} \cdot \sqrt{S_E^2 + (\Delta x)^2} \label{eq:fehler_messreihe_einfach}
\end{equation}

\section{Güte eines Fits}
Um die Güte eines Fits zu beurteilen, bildet man häufig die sogenannte Chi-Quadrat-Summe $\chi^2$, welche die Abweichung der Messwerte von einer Fitfunktion beschreibt. Das Vorgehen ist ähnlich zur Bestimmung der $\sigma$-Abweichung zweier Werte. Seien $x_i \in X$ mit $X$ eine Menge von Daten, bezüglich derer Messwerte augenommen wurden. Sei $f:~X \rightarrow Y$ eine Fitfunktion und $g:~X \rightarrow Y$ die Funktion, welche jedem Wert $x_i$ den entsprechenden Messwert $g(x_i)$ zuordnet. $Y$ kann bspw. $\mathbb{R}$ oder $\mathbb{C}$ darstellen. Sei $\Delta g(x_i)$ der Fehler zu $g(x_i)$. Dann gilt:
\begin{equation}
    \chi^2 = \sum_{x_i \in X} \biggl(\frac{f(x_i) - g(x_i)}{\Delta g(x_i)}\biggr)^2~\label{eq:chisquare}
\end{equation}
Diese sollte möglichst klein sein, damit die Güte des Fits hoch ist.

Noch besser für eine Diskussion der Güte des Fits geeignet ist die reduzierte Chi-Quadrat-Summe $\chi^2_{red}$, welche die Chi-Quadrat-Summe anhand der Anzahl der Messwerte normiert:
\begin{equation}
    \chi^2_{red} = \frac{\chi^2}{\text{dof}}~\label{eq:chisquare_red}
\end{equation}
Man möchte einen Wert $\chi^2_{red} \approx 1$. Grundsätzlich bedeutet $\chi^2_{red} < 1$ ein „overfitting”, welches bei zu viel Rauschen entstehen kann und $\chi^2_{red} >> 1$ bedeutet einen allgemein schlechten Fit.

Letztlich kann noch die „Fitwahrscheinlichkeit” $P$ berechnet werden, welche die Wahrscheinlichkeit beschreibt, dass bei einer Wiederholungsmessung einer gleicher oder größerer $\chi^2$-Wert bestimmt wird. Diese wird mit Hilfe der Funktion chi2 aus dem Python-Paket scipy.stats bestimmt:
\begin{equation}
    P = (1 - chi2(\chi^2, \text{dof})) \cdot 100~\label{eq:fitwahrscheinkichkeit}
\end{equation} 